% SC26 MLIR tutorial -- Session 2, Part 2 (distributing Linalg on the GPU).
% Presenter deck.
%
% Build:  make            (the deck PDF)
% See presentation/README.md.
\providecommand{\notesmode}{0}
\documentclass[aspectratio=169,10pt]{beamer}
\usepackage{sc26}
\ifnum\notesmode=1 \setbeameroption{show notes}\fi

\title{Introduction to MLIR}
\subtitle{Session 2, Part 2 --- Distributing Linalg on the GPU}
\author{SC26 Tutorial}
\date{}

\begin{document}

\begin{frame}[plain]
  \titlepage
\end{frame}
\note{Pre-run setup cells 0--9 of the walkthrough notebook on a T4 runtime
before the room fills: the toolchain download is the slowest thing in this
part and teaches nothing.

One exercise, widening the distribution schedule: about 10 minutes of
keyboard time. Step 1 is read and run together; nobody types anything in it.

Known problem: the walkthrough notebook clones session2-part2-solution, whose
schedule is for matmul only. Its matvec cell (19) stops on the race check
("tiling is not thread safe at axis \#1") on the first pass. Until that is
settled, run step 2's matvec and provided-schedule matmul cells (19 to 24)
from the attendee notebook.}

% =======================================================================
\begin{frame}{Setup}
  \begin{itemize}
    \item<1-> \textbf{Switch the runtime to a GPU:}
      \cd{Change runtime type $\to$ T4 GPU}, before running any cell.
    \item<2-> \textbf{It is a fresh runtime.} Nothing from Part 1 carries
      over, so the notebook downloads the toolchain again.
    \item<3-> \textbf{You edit one MLIR file, a schedule,} as in Part 1: edit,
      save, re-run the cell. No C++.
    \item<4-> \textbf{One exercise:} distribute matmul along its columns too.
  \end{itemize}
\end{frame}
\note{Change runtime type is under the Runtime menu, or the arrow next to
Connect at the top right. Ctrl+F8 (run every cell above) catches someone up
before the exercise; after it, not on the exercise's cells, for the reason in
the exercise slide's note.}

% =======================================================================
\begin{frame}{Same Tiling, onto the GPU}
  \begin{center}
  \begin{tikzpicture}[
      box/.style={draw=scblue,fill=scwash,rounded corners=2pt,inner sep=5pt,
                  font=\ttfamily\scriptsize,minimum height=2.4em,align=center},
      lbl/.style={font=\scriptsize\itshape,color=scgrey,align=center},
      arr/.style={-{Stealth[length=2mm]},draw=scgrey,thick}]

    \node[box,draw=scaccent] (op) at (0,0) {linalg.matmul};
    \node[box] (for) at (3.9,1.0) {scf.for\\\normalfont\itshape tile loops};
    \node[box] (cpu) at (11.0,1.0) {JIT + run\\\normalfont\itshape CPU, one thread};
    \draw[arr] (op) -- (for);
    \node[lbl] at (1.35,0.95) {tile\_using\_for\\Part 1};
    \draw[arr] (for) -- node[lbl,above]{loops + CPU pipeline} (cpu);

    \visible<2->{
      \node[box,draw=scaccent] (forall) at (3.9,-1.0)
        {scf.forall\\\normalfont\itshape tile loops};
      \node[box] (launch) at (7.5,-1.0)
        {gpu.launch\\\normalfont\itshape blocks, threads};
      \node[box] (t4) at (11.0,-1.0) {JIT + run\\\normalfont\itshape T4 GPU};
      \draw[arr,draw=scaccent] (op) -- (forall);
      \node[lbl] at (1.35,-0.95) {tile\_using\_forall\\\emph{this part}};
      \draw[arr] (forall) -- (launch);
      \draw[arr] (launch) -- (t4);
    }
  \end{tikzpicture}
  \end{center}
  \vspace{0.3em}
  \begin{itemize}
    \item<1-> In Part 1, we tiled Linalg ops into \cd{scf.for} loops: one tile after
      another, on one CPU thread.
    \item<2-> Here we tile the same ops into \cd{scf.forall} loops, whose
      iterations may run in parallel.
    \item<3-> The schedule also says which loops become the blocks and which
      the threads, and builds the \cd{gpu.launch}.
  \end{itemize}
\end{frame}
\note{Step 1 is the payloads: Part 1's two files with the host registration
from Session 1 Part 2. Step 2 is the schedule, including how it maps the
loops onto the GPU, the exercise, and why the mapping it asks for is the
right one for memory.}

% =======================================================================
\section{Step 1 --- The GPU Payloads}
\begin{frame}[plain]\sectionpage\end{frame}

\begin{frame}[fragile]{The GPU Payloads}
  \cd{session2/linalg\_matvec\_gpu.mlir}:
\begin{lstlisting}[language=mlir]
module {
  func.func @matvec(%mat: memref<128x512xf32>, %vec: memref<512xf32>,
                    %out: memref<128xf32>) {
    (@\hl<2>{\%mat\_u = \lkw{memref.cast} \%mat : \lkw{memref}<128x512xf32> \lkw{to} \lkw{memref}<*xf32>}@)
    (@\hl<2>{\%vec\_u = \lkw{memref.cast} \%vec : \lkw{memref}<512xf32> \lkw{to} \lkw{memref}<*xf32>}@)
    (@\hl<2>{\%out\_u = \lkw{memref.cast} \%out : \lkw{memref}<128xf32> \lkw{to} \lkw{memref}<*xf32>}@)
    (@\hl<2>{\lkw{gpu.host\_register} \%mat\_u : \lkw{memref}<*xf32>}@)
    (@\hl<2>{\lkw{gpu.host\_register} \%vec\_u : \lkw{memref}<*xf32>}@)
    (@\hl<2>{\lkw{gpu.host\_register} \%out\_u : \lkw{memref}<*xf32>}@)
    (@\hl<3>{\lkw{linalg.matvec} ins(\%mat, \%vec : \lkw{memref}<128x512xf32>, \lkw{memref}<512xf32>)}@)
                  (@\hl<3>{outs(\%out : \lkw{memref}<128xf32>)}@)
    return
  }
}
\end{lstlisting}
  \vspace{0.2em}
  \begin{itemize}
    \item<2-> The host buffers, registered so the GPU can reach them, as in
      Session 1 Part 2's kernel. Boilerplate: we leave it there.
    \item<3-> The op is Part 1's \cd{linalg.matvec}, unchanged.
  \end{itemize}
\end{frame}
\note{As in the repo, without the header comment. linalg\_matmul\_gpu.mlir is
the same, around Part 1's linalg.matmul. Both files are complete; round-trip
them in the notebook (cell 13) as you show them. The cast to an unranked
memref is only because gpu.host\_register takes one.

Nothing unregisters a buffer. If a registered buffer is freed and a later
numpy array lands on its address, the kernel writes somewhere numpy is not
reading and the result comes back all zeros. That is why the notebook keeps
every buffer it hands a kernel in REGISTERED.}

% =======================================================================
\section{Step 2 --- Distributing, with a Transform Schedule}
\begin{frame}[plain]\sectionpage\end{frame}

\begin{frame}[fragile]{Tiling into Parallel Loops}
  \begin{center}
    Tiling matvec's \cd{i} by 32, the shape of it:
  \end{center}
  \begin{columns}[t]
    \begin{column}{0.47\textwidth}
      \begin{center}\small\textbf{Part 1: \cd{tile\_using\_for}}\end{center}
\begin{lstlisting}[language=mlir]
(@\hl<2>{\lkw{scf.for} \%i = 0 \lkw{to} 128 \lkw{step} 32 \{}@)
  %o = memref.subview %out[%i] [32]
  linalg.matvec ins(...) outs(%o)
}
\end{lstlisting}
    \end{column}
    \begin{column}{0.47\textwidth}
      \begin{center}\small\textbf{this part: \cd{tile\_using\_forall}}\end{center}
\begin{lstlisting}[language=mlir]
(@\hl<3>{\lkw{scf.forall} (\%b) in (4) \{}@)
  (@\hl<3>{\%i = \lkw{affine.apply} \%b * 32}@)
  %o = memref.subview %out[%i] [32]
  linalg.matvec ins(...) outs(%o)
} (@\hl<4>{\{mapping = [\#gpu.block<x>]\}}@)
\end{lstlisting}
    \end{column}
  \end{columns}
  \vspace{0.6em}
  \begin{itemize}
    \item<2-> \cd{scf.for} runs its iterations one after another, in order.
    \item<3-> \cd{scf.forall} says its iterations are independent: they may
      run in any order, or all at once.
    \item<4-> Its \alert{mapping} says what runs each iteration on the GPU:
      here a block, numbered by its block id in x.
    \item<5-> Only a \cd{parallel} loop can be tiled this way: tiles of a
      reduction would add into the same \cd{out[i]} at once.
      \cd{tile\_using\_forall} only warns: \cd{tiling is not thread safe}.
  \end{itemize}
\end{frame}
\note{Both listings are the shape, not real output: the names are made up and
the types left out. The mapping attribute is printed exactly like that on a
real scf.forall.

The forall counts tiles, 4, and each iteration computes its own offset; the
for loop steps through the offsets directly.

The real warning is "warning: tiling is not thread safe at axis \#1" (the
axis is the loop's position). mlir-opt still emits the loops and exits 0,
which is why every notebook cell that applies the schedule keeps its stderr
and fails on that text.}

% =======================================================================
\begin{frame}{Blocks and Threads, from Two Tilings}
  \begin{center}
  \begin{tikzpicture}[
      lbl/.style={font=\scriptsize,color=scink},
      ids/.style={font=\ttfamily\scriptsize,color=scaccent}]
    % out, 128 elements: 4 blocks of 32 rows, top to bottom
    \foreach \b in {0,...,3} {
      \draw[scblue,fill=scwash] (0,1.6-0.8*\b) rectangle (0.45,0.8-0.8*\b);
      \node[ids,anchor=east] at (-0.1,1.2-0.8*\b) {bx = \b};
    }
    \node[ids] at (0.225,1.95) {\#gpu.block<x>};
    \node[lbl] at (0.225,-1.9) {\cd{out}, 128};

    \visible<2->{
      % block 1, zoomed: 32 threads, one row each
      \fill[scaccent!25] (0,0.8) rectangle (0.45,0);
      \draw[scblue] (0,0.8) rectangle (0.45,0);
      \fill[scaccent!25] (2.6,1.6) rectangle (3.05,-1.6);
      \foreach \t in {1,...,31} {
        \draw[scblue!40,very thin] (2.6,1.6-0.1*\t) -- (3.05,1.6-0.1*\t);
      }
      \draw[scblue] (2.6,1.6) rectangle (3.05,-1.6);
      \draw[scgrey,dashed] (0.45,0.8) -- (2.6,1.6);
      \draw[scgrey,dashed] (0.45,0) -- (2.6,-1.6);
      \node[ids] at (2.825,1.95) {\#gpu.thread<x>};
      \node[ids,anchor=west] at (3.12,1.55) {tx = 0};
      \node[ids,anchor=west] at (3.12,-1.55) {tx = 31};
      \node[lbl] at (2.825,-1.9) {block 1: 32 threads};
    }

    \visible<3->{
      % one thread: tx = 5 of block 1
      \fill[scaccent] (2.6,1.1) rectangle (3.05,1.0);
      \node[ids,anchor=east] at (2.53,1.05) {tx = 5};
      \draw[-{Stealth[length=1.5mm]},scaccent] (3.15,1.05) -- (4.45,1.05);
      \node[lbl,anchor=west,align=left] at (4.55,1.05)
        {one thread: \cd{out[i] += mat[i, j] * vec[j]},\\for all 512 \cd{j}};
      \node[lbl,anchor=west] at (4.55,0.2)
        {\cd{i = bx * 32 + tx = 1 * 32 + 5 = 37}};
    }
  \end{tikzpicture}
  \end{center}
  \vspace{0.4em}
  \begin{itemize}
    \item<1-> Tile \cd{i} by 32, mapped to \cd{\#gpu.block<x>}: 4 iterations,
      the grid's 4 blocks, \cd{bx} = 0 \dots{} 3.
    \item<2-> Tile each of those by 1, mapped to \cd{\#gpu.thread<x>}: 32
      iterations, the 32 threads of each block, \cd{tx} = 0 \dots{} 31.
    \item<3-> Each thread computes one \cd{out[i]}, reducing over \cd{j}
      itself. Its \cd{i} is Session 1 Part 2's
      \cd{blockIdx * blockDim + threadIdx}.
  \end{itemize}
\end{frame}
\note{Recall from Session 1 Part 2: a launch is a grid of blocks, each block a
group of threads, and every thread knows its block id, its thread id and the
block size.

4 blocks of 32 threads is 128 threads, one per element of out: the same work
as Session 1 Part 2's matvec kernel, which used one block of numRows threads.

Not to scale: out is drawn as four boxes, and the zoomed block has its 32
rows drawn large.}

% =======================================================================
\begin{frame}[fragile]{Reference: \cd{distribute\_linalg.mlir}, Step 1 of 4 --- The Tilings}
\begin{lstlisting}[language=mlir]
module attributes {transform.with_named_sequence} {
  transform.named_sequence @__transform_main(
      %root: !transform.any_op {transform.readonly}) {
    (@\hl<2>{\%ops = \lkw{transform.structured.match} interface\{LinalgOp\} in \%root}@)
      : (!transform.any_op) -> !transform.any_op
    (@\hl<5>{\%block\_op}@), %blocks = (@\hl<3>{\lkw{transform.structured.tile\_using\_forall} \%ops}@)
      (@\hl<3>{tile\_sizes [32] (mapping = [\#gpu.block<x>])}@)
      : (!transform.any_op) -> (!transform.any_op, !transform.any_op)
    %thread_op, %threads = (@\hl<4>{\lkw{transform.structured.tile\_using\_forall} \hl<5>{\%block\_op}}@)
      (@\hl<4>{tile\_sizes [1] (mapping = [\#gpu.thread<x>])}@)
      : (!transform.any_op) -> (!transform.any_op, !transform.any_op)
    // next: build the kernel
    transform.yield
  }
}
\end{lstlisting}
  \vspace{0.3em}
  \begin{itemize}
    \item<2-> \cd{match} by interface, as in Part 1: any Linalg op.
    \item<3-> The blocks: tile \cd{i} by 32 into an \cd{scf.forall}. The
      mapping has one entry per tiled loop: this one becomes the blocks, by
      block id in x.
    \item<4-> The threads: the op inside each block, by 1. Each iteration
      becomes a thread of that block, by thread id in x.
    \item<5-> \cd{\%block\_op} is the op inside the blocks: the first step
      returns it, and the second tiles it.
  \end{itemize}
\end{frame}
\note{session2/distribute\_linalg.mlir, as in the repo without its comments;
the "next" comment stands for the lines on the next slide. It runs the way
Part 1's schedule did: -{}-transform-preload-library, then
-{}-transform-interpreter (cell 19). Like Part 1's, it is a working
schedule, written for matvec; they read it before changing anything.

\%blocks and \%threads are handles to the two scf.foralls; nothing uses them.}

% =======================================================================
\begin{frame}[fragile]{Reference: \cd{distribute\_linalg.mlir}, Step 2 of 4 --- Building the Kernel}
\begin{lstlisting}[language=mlir]
module attributes {transform.with_named_sequence} {
  transform.named_sequence @__transform_main(...) {
    // the match, then the block and thread tilings
    (@\hl<2>{\%func = \lkw{transform.structured.match} ops\{["func.func"]\} in \%root}@)
      : (!transform.any_op) -> !transform.any_op
    %launch = (@\hl<3>{\lkw{transform.gpu.map\_forall\_to\_blocks} \%func generate\_gpu\_launch}@)
      : (!transform.any_op) -> !transform.any_op
    (@\hl<4>{\lkw{transform.gpu.map\_nested\_forall\_to\_threads} \%launch block\_dims = [32, 1, 1]}@)
      sync_after_distribute = false
      : (!transform.any_op) -> !transform.any_op
    transform.yield
  }
}
\end{lstlisting}
  \vspace{0.3em}
  \begin{itemize}
    \item<2-> Match the function: the next two steps work on the code around
      the loops.
    \item<3-> \cd{map\_forall\_to\_blocks} puts the block loop into a
      \cd{gpu.launch}, one block per iteration, and replaces its index with
      \cd{gpu.block\_id x}.
    \item<4-> \cd{map\_nested\_forall\_to\_threads} does the same for the
      thread loop, with \cd{gpu.thread\_id x}. \cd{block\_dims} is the size of
      each block: 32 threads in x.
  \end{itemize}
\end{frame}
\note{The rest of the same file. sync\_after\_distribute = false leaves out a
gpu.barrier after the thread loop, which this kernel does not need: no thread
reads what another wrote.

Both ops read the mapping attributes the tilings put on the loops. A loop
without one is an error: "scf.forall op requires a mapping attribute".}

% =======================================================================
\begin{frame}[fragile]{Reference: \cd{distribute\_linalg.mlir}, Step 3 of 4 --- Applied to Matvec}
\begin{lstlisting}[language=mlir,basicstyle=\ttfamily\fontsize{6.6}{7.9}\selectfont]
(@\hl<2>{\lkw{gpu.launch} blocks(\%arg3, \%arg4, \%arg5) in (\%arg9 = \%c4, \%arg10 = \%c1\_4, \%arg11 = \%c1\_5)}@)
           (@\hl<2>{threads(\%arg6, \%arg7, \%arg8) in (\%arg12 = \%c32, \%arg13 = \%c1\_2, \%arg14 = \%c1\_3) \{}@)
  (@\hl<3>{\%block\_id\_x = \lkw{gpu.block\_id}  x}@)
  ...
  (@\hl<3>{\%0 = \lkw{affine.apply} \lkw{affine\_map}<(d0) -> (d0 * 32)>(\%block\_id\_x)}@)
  %subview = memref.subview %arg0[(@\hl<3>{\%0}@), 0] [32, 512] [1, 1] : ...
  %subview_7 = memref.subview %arg1[0] [512] [1] : ...
  %subview_8 = memref.subview %arg2[(@\hl<3>{\%0}@)] [32] [1] : ...
  (@\hl<4>{\%thread\_id\_x = \lkw{gpu.thread\_id}  x}@)
  ...
  %subview_9 = memref.subview %subview[(@\hl<4>{\%thread\_id\_x}@), 0] [1, 512] [1, 1] : ...
  %subview_10 = memref.subview %subview_7[0] [512] [1] : ...
  %subview_11 = memref.subview %subview_8[(@\hl<4>{\%thread\_id\_x}@)] [1] [1] : ...
  (@\hl<5>{\lkw{linalg.matvec} ins(\%subview\_9, \%subview\_10 : ...) outs(\%subview\_11 : ...)}@)
  gpu.terminator
}
\end{lstlisting}
  \vspace{0.3em}
  \begin{itemize}
    \item<2-> Grid size: \cd{(4, 1, 1)}, block size: \cd{(32, 1, 1)}: the
      block loop's 4 iterations, and \cd{block\_dims}.
    \item<3-> Block \cd{bx} (\cd{\%block\_id\_x}): 32 rows of \cd{mat} and
      \cd{out}, from row \cd{\%0 = bx * 32}.
    \item<4-> Thread \cd{tx} (\cd{\%thread\_id\_x}): one row of that slice,
      row \cd{bx * 32 + tx}.
    \item<5-> Inside, the original \cd{linalg.matvec}, on that one row. Still
      Linalg: the next step lowers it.
  \end{itemize}
\end{frame}
\note{The real output of cell 19, with mlir-opt -{}-mlir-print-local-scope so
the affine map is written out in place; the notebook prints it as \#map,
defined at the top. Compressed: each "..." after a colon is a type, the bare
"..." lines are the unused y and z ids and two constants, and the casts and
host registrations above the launch are left out.

vec is used whole because its map, (i, j) -> (j), has no i in it: tiling i
does not slice it.}

% =======================================================================
\begin{frame}[fragile]{Reference: Step 4 of 4 --- Compiling and Execution}
\begin{lstlisting}[style=plain,basicstyle=\ttfamily\scriptsize]
(@\hl<2>{convert-linalg-to-loops}@)
(@\hl<3>{expand-strided-metadata}@)
(@\hl<3>{lower-affine}@)
(@\hl<4>{gpu-lower-to-nvvm-pipeline\{cubin-chip=sm\_75 cubin-format=bin cubin-features=+ptx78\}}@)
\end{lstlisting}
  \vspace{0.3em}
  \begin{itemize}
    \item<2-> \cd{convert-linalg-to-loops} turns the \cd{linalg.matvec} into
      the thread's loop over \cd{j}: Session 1 Part 2's matvec kernel,
      generated.
    \item<3-> Part 1's two passes, for the subviews and their
      \cd{affine.apply}.
    \item<4-> Session 1 Part 2's GPU pipeline: the body compiled for this GPU,
      the host side rewritten to launch it.
    \item<5-> JIT-compiled as in Session 1 Part 2 --- the CUDA runtime library
      in \cd{shared\_libs}, then \cd{engine.initialize()} --- and checked
      against numpy:
  \end{itemize}
  \vspace{0.2em}
  \visible<5->{\begin{center}\color{scgreen}\large\ttfamily OK\end{center}}
\end{frame}
\note{Cell 21. GPU\_LOWERING\_PIPELINE is these passes, with
llvm-request-c-wrappers as well. The notebook reads the chip from nvidia-smi
rather than hard-coding sm\_75.

Lowered, the body is a one-trip loop over the thread's row around the
512-step reduction over j.

Without engine.initialize() every CUDA call fails with
CUDA\_ERROR\_INVALID\_HANDLE: the compiled kernel is loaded by a global
constructor that nothing else runs.}

% =======================================================================
\begin{frame}[fragile]{The Same Schedule, on Matmul}
  \begin{center}
  \begin{tikzpicture}[
      lbl/.style={font=\scriptsize,color=scink},
      ids/.style={font=\ttfamily\scriptsize,color=scaccent}]
    % c, 128 x 256, as 4 x 8 tiles of 32 x 32. Left: the provided schedule.
    \fill[scaccent!25] (0,0.45) rectangle (3.6,0);
    \fill[scaccent] (0,0.31) rectangle (3.6,0.25);
    \draw[step=0.45,scgrey!45,very thin] (0,-0.9) grid (3.6,0.9);
    \foreach \b in {0,...,3} {
      \draw[scblue] (0,0.9-0.45*\b) rectangle (3.6,0.45-0.45*\b);
      \node[ids,anchor=east] at (-0.1,0.675-0.45*\b) {bx = \b};
    }
    \node[lbl] at (1.8,1.25) {provided: 4 blocks of 32 threads};
    \node[lbl,align=center] at (1.8,-1.35)
      {one thread: a whole row of \cd{c}};

    \visible<3->{
      % right: the exercise. A grid of 8 x 4 blocks, one thread per element
      \begin{scope}[xshift=6.6cm]
        \fill[scaccent!25] (0.9,0.45) rectangle (1.35,0);
        \fill[scaccent] (1.0,0.31) rectangle (1.06,0.25);
        \draw[step=0.45,scblue] (0,-0.9) grid (3.6,0.9);
        \draw[-{Stealth[length=1.2mm]},scaccent] (0,1.0) -- (3.6,1.0)
          node[right,ids]{x};
        \draw[-{Stealth[length=1.2mm]},scaccent] (-0.12,0.9) -- (-0.12,-0.9)
          node[below,ids]{y};
        \node[lbl] at (1.8,1.35) {the exercise: blocks of 32 $\times$ 32
          threads};
        \node[lbl,align=center] at (1.8,-1.35)
          {one thread: one \cd{c[i, j]}};
      \end{scope}
    }
  \end{tikzpicture}
  \end{center}
  \vspace{0.1em}
  The innermost op, from the same schedule:
\begin{lstlisting}[language=mlir,basicstyle=\ttfamily\fontsize{6.6}{7.9}\selectfont]
linalg.matmul ins(%subview_9, %subview_10 : memref<1x512xf32, ...>, memref<512x256xf32, ...>)
              outs(%subview_11 : memref<(@\hl<2>{1x256}@)xf32, ...>)
\end{lstlisting}
  \vspace{0.1em}
  \begin{itemize}
    \item<1-> Unchanged, the schedule tiles only the first loop, \cd{i}, and
      maps it to x: 4 blocks of 32 threads, as for matvec.
    \item<2-> Each thread gets a 1 $\times$ \alert{256} slice of \cd{c}: a
      whole row, with \cd{j} and \cd{k} as loops in its body.
    \item<3-> The exercise: distribute \cd{j} too, so that each thread
      computes one \cd{c[i, j]}, with the columns on x and the rows on y.
  \end{itemize}
\end{frame}
\note{Cells 23 and 24: the launch, and a run that passes the numpy check. It
is correct, just 128 threads for 32 768 elements, each doing 256 x 512
multiply-adds.

The grids are to scale, one cell per 32 by 32 tile of c; the thread's row and
element are drawn larger than they are.}

% =======================================================================
\begin{frame}[fragile]{Exercise 7 --- Distribute Matmul along Columns Too}
  \begin{center}\large
    Modify \cd{session2/distribute\_linalg.mlir} so that each thread computes
    one \mbox{\cd{c[i, j]}}, in blocks of 32 $\times$ 32 threads: the
    columns (\cd{j}) on x, the rows (\cd{i}) on y.
  \end{center}
  \vspace{0.1em}
  \textbf{What is different from matvec:}
  \begin{itemize}
    \item Matvec has one parallel loop, \cd{i}. Matmul has two, \cd{i} and
      \cd{j}; \cd{k} is a reduction, and stays inside each thread.
    \item One tile size and one mapping entry per loop, in
      \cd{iterator\_types} order: \cd{i} first, then \cd{j}.
    \item Each block is 32 $\times$ 32 threads.
  \end{itemize}
  \vspace{0.3em}
  \textbf{Check as you go --- round-trip it} after every save:
\begin{lstlisting}[style=plain,basicstyle=\ttfamily\small]
!mlir-opt distribute_linalg.mlir
\end{lstlisting}
  \vspace{0.1em}
  \begin{itemize}
    \item That checks the schedule is well-formed, not that it does what you
      meant. For that, run the check under \emph{Matmul, distributed along
      columns too}: matmul's innermost op must write a 1 $\times$ 1 slice of
      \cd{c}, cut at \cd{[\%thread\_id\_y, \%thread\_id\_x]}.
  \end{itemize}
  \vspace{0.3em}
  \begin{center}\color{scaccent}\large
    10 minutes to write it --- then we walk through the solution.
  \end{center}
\end{frame}
\note{After an edit: the four cells under "Matmul, distributed along columns
too", in order: the round-trip (cell 26), the apply (28), the check (30) and
the run (32). Not Ctrl+F8 on any of them: it re-runs the matvec cell (19),
and a schedule widened for matmul distributes matvec's reduction there, so
that cell stops on "tiling is not thread safe at axis \#1" and the race
message.

What the usual mistakes produce:
block\_dims left at [32, 1, 1]: "the number of required parallel resources
(blocks or threads) 1024 overflows the number of available resources 32".
A second size with no second mapping entry: "scf.forall op requires a
mapping attribute". x and y the wrong way round: the check stops with "x is
not on the columns yet", though it would run and match numpy. A separate step for
j, as in Part 1: "scf.forall op requires a mapping attribute of kind
'thread' or 'warp'". [32, 32, 32]: "tiling is not thread safe at axis \#2".

Until the exercise is done, the check stops with "matmul is not distributed
along columns yet -- each thread still computes a 1x256 slice of c".}

% =======================================================================
\begin{frame}[fragile]{Solution --- \cd{distribute\_linalg.mlir}, and What Changed}
\begin{lstlisting}[language=mlir,basicstyle=\ttfamily\fontsize{6.4}{7.6}\selectfont]
module attributes {transform.with_named_sequence} {
  transform.named_sequence @__transform_main(
      %root: !transform.any_op {transform.readonly}) {
    %ops = transform.structured.match interface{LinalgOp} in %root
      : (!transform.any_op) -> !transform.any_op
    %block_op, %blocks = transform.structured.tile_using_forall %ops
      tile_sizes (@\df{[32, 32]}@) (mapping = (@\df{[\#gpu.block<y>, \#gpu.block<x>]}@))
      : (!transform.any_op) -> (!transform.any_op, !transform.any_op)
    %thread_op, %threads = transform.structured.tile_using_forall %block_op
      tile_sizes (@\df{[1, 1]}@) (mapping = (@\df{[\#gpu.thread<y>, \#gpu.thread<x>]}@))
      : (!transform.any_op) -> (!transform.any_op, !transform.any_op)
    %func = transform.structured.match ops{["func.func"]} in %root
      : (!transform.any_op) -> !transform.any_op
    %launch = transform.gpu.map_forall_to_blocks %func generate_gpu_launch
      : (!transform.any_op) -> !transform.any_op
    transform.gpu.map_nested_forall_to_threads %launch block_dims = (@\df{[32, 32, 1]}@)
      sync_after_distribute = false
      : (!transform.any_op) -> !transform.any_op
    transform.yield
  }
}
\end{lstlisting}
  \vspace{0.2em}
  \begin{itemize}
    \item<2-> Both tilings take \cd{j} as well: \cd{[32, 32]}, then
      \cd{[1, 1]}.
    \item<3-> The mappings follow the loops: \cd{i} first, on y; \cd{j}
      second, on x.
    \item<4-> Each block is 32 $\times$ 32 threads.
  \end{itemize}
\end{frame}
\note{The handles do not change: tile\_using\_forall makes one scf.forall
however many loops it tiles.

This schedule is for matmul only. On matvec its second size lands on the
reduction, and the matvec cell stops on the race check.}

% =======================================================================
\begin{frame}[fragile]{Solution --- \cd{gpu.launch}}
\begin{lstlisting}[language=mlir,basicstyle=\ttfamily\fontsize{6.1}{7.0}\selectfont]
(@\hl<2>{\lkw{gpu.launch} blocks(\%arg3, \%arg4, \%arg5) in (\%arg9 = \%c8, \%arg10 = \%c4, \%arg11 = \%c1\_4)}@)
           (@\hl<2>{threads(\%arg6, \%arg7, \%arg8) in (\%arg12 = \%c32, \%arg13 = \%c32\_2, \%arg14 = \%c1\_3) \{}@)
  (@\hl<3>{\%block\_id\_x = \lkw{gpu.block\_id}  x}@)
  (@\hl<3>{\%block\_id\_y = \lkw{gpu.block\_id}  y}@)
  ...
  (@\hl<3>{\%0 = \lkw{affine.apply} \lkw{affine\_map}<(d0) -> (d0 * 32)>(\%block\_id\_y)}@)
  (@\hl<3>{\%1 = \lkw{affine.apply} \lkw{affine\_map}<(d0) -> (d0 * 32)>(\%block\_id\_x)}@)
  %subview = memref.subview %arg0[(@\hl<3>{\%0}@), 0] [32, 512] [1, 1] : ...
  %subview_6 = memref.subview %arg1[0, (@\hl<3>{\%1}@)] [512, 32] [1, 1] : ...
  %subview_7 = memref.subview %arg2[(@\hl<3>{\%0, \%1}@)] [32, 32] [1, 1] : ...
  (@\hl<4>{\%thread\_id\_x = \lkw{gpu.thread\_id}  x}@)
  (@\hl<4>{\%thread\_id\_y = \lkw{gpu.thread\_id}  y}@)
  ...
  %subview_8 = memref.subview %subview[(@\hl<4>{\%thread\_id\_y}@), 0] [1, 512] [1, 1] : ...
  %subview_9 = memref.subview %subview_6[0, (@\hl<4>{\%thread\_id\_x}@)] [512, 1] [1, 1] : ...
  %subview_10 = memref.subview %subview_7[(@\hl<4>{\%thread\_id\_y, \%thread\_id\_x}@)] [1, 1] [1, 1] : ...
  (@\hl<5>{\lkw{linalg.matmul} ins(\%subview\_8, \%subview\_9 : ...) outs(\%subview\_10 : ...)}@)
  gpu.terminator
}
\end{lstlisting}
  \begin{itemize}
    \item<2-> Grid size: \cd{(8, 4, 1)}, block size: \cd{(32, 32, 1)}: x runs
      over the columns.
    \item<3-> Block \cd{(by, bx)} computes the 32 $\times$ 32 tile of \cd{c}
      at \cd{(by * 32, bx * 32)}, reading 32 rows of \cd{a} and 32 columns
      of \cd{b}.
    \item<4-> Thread \cd{(ty, tx)} computes element \cd{[ty, tx]} of that
      tile, reading one row of \cd{a} and one column of \cd{b}.
    \item<5-> Lowered, that is 512 iterations of \cd{k}. Run, and checked
      against numpy: \visible<5->{{\color{scgreen}\ttfamily OK}}
  \end{itemize}
\end{frame}
\note{The real output of the apply cell (28) with the solved schedule, with
the affine maps written out in place (-{}-mlir-print-local-scope) and
compressed as for matvec. Grid (8, 4, 1): x counts the 8 blocks of columns,
y the 4 blocks of rows. by cuts the 32-row strips of a and c, bx the
32-column strips of b and c; the thread ids then take one row of a's strip
and one column of b's. Cell 30 checks it and cell 32 runs it. Run cells 26
to 32 in the walkthrough notebook as you go.}

% =======================================================================
\begin{frame}{Which Thread Computes Which Element}
  \begin{center}
  \begin{tikzpicture}[
      cell/.style={draw=scblue,fill=scwash,minimum width=1.05cm,
                   minimum height=0.5cm,font=\ttfamily\scriptsize},
      hdr/.style={font=\ttfamily\scriptsize,color=scaccent}]
    \visible<2->{
      \fill[scaccent!30,rounded corners=2pt] (-0.62,0.32) rectangle (6.25,-0.32);
    }
    \foreach \c in {0,...,4} {
      \node[hdr] at (\c*1.15, 0.6) {tx = \c};
      \foreach \r in {0,...,4} {
        \node[cell] at (\c*1.15, -\r*0.6) {(\r, \c)};
      }
    }
    \foreach \r in {0,...,4} {
      \node[hdr,anchor=east] at (-0.7, -\r*0.6) {ty = \r};
      \node[font=\scriptsize,color=scgrey] at (5.9, -\r*0.6) {$\cdots$};
    }
    \foreach \c in {0,...,4} {
      \node[font=\scriptsize,color=scgrey] at (\c*1.15, -2.95) {$\vdots$};
    }
    \node[font=\scriptsize,color=scgrey,anchor=west] at (6.6, -0.95)
      {a block's 32 $\times$ 32 tile of \cd{c}:};
    \node[font=\scriptsize,color=scgrey,anchor=west] at (6.6, -1.3)
      {each cell is the element \cd{(i, j)}};
    \visible<2->{
      \node[font=\scriptsize,color=scaccent,anchor=west] at (6.6, 0)
        {one warp: \cd{ty} = 0, \cd{tx} = 0 \dots{} 31};
    }
  \end{tikzpicture}
  \end{center}
  \vspace{0.4em}
  \begin{itemize}
    \item<1-> The y ids pick the row of \cd{c}, \cd{by * 32 + ty}; the x ids
      pick the \alert{column}, \cd{bx * 32 + tx}. The schedule said so.
    \item<2-> A warp is 32 threads with consecutive \cd{tx} and the same
      \cd{ty}: here, 32 elements along a row of \cd{c}.
  \end{itemize}
\end{frame}
\note{The same picture as Session 1 Part 2's "Thread to Element Mapping",
block by block: that kernel also put x on the columns.}

% =======================================================================
\begin{frame}{Why x on the Columns: Coalescing}
  \begin{center}
  \begin{tikzpicture}[
      lbl/.style={font=\scriptsize,color=scink,align=center}]
    % a corner of c, row-major. Left: a warp with x on the columns.
    \fill[scaccent!45] (0,-0.3) rectangle (3.0,-0.6);
    \draw[step=0.3,scblue!50,very thin] (0,-1.8) grid (3.0,0);
    \draw[scblue] (0,-1.8) rectangle (3.0,0);
    \draw[-{Stealth[length=1.2mm]},scgrey] (0,0.2) -- (3.0,0.2)
      node[right,font=\scriptsize,color=scgrey]{memory order};
    \node[lbl] at (1.5,0.6) {x on the columns: this schedule};
    \node[lbl] at (1.5,-2.25) {a warp: 32 adjacent floats, 128 bytes};

    \visible<3->{
      % right: a warp with x on the rows
      \begin{scope}[xshift=5.8cm]
        \fill[scaccent!45] (0.6,0) rectangle (0.9,-1.8);
        \draw[step=0.3,scblue!50,very thin] (0,-1.8) grid (3.0,0);
        \draw[scblue] (0,-1.8) rectangle (3.0,0);
        \draw[-{Stealth[length=1.2mm]},scgrey] (0,0.2) -- (3.0,0.2);
        \node[lbl] at (1.5,0.6) {x on the rows: y and x swapped};
        \node[lbl] at (1.5,-2.25) {a warp: 32 floats, a row (1 KB) apart};
      \end{scope}
    }
  \end{tikzpicture}
  \end{center}
  \vspace{0.3em}
  \begin{itemize}
    \item<1-> \cd{c} is row-major: \cd{c[i, j]} and \cd{c[i, j + 1]} are
      adjacent in memory.
    \item<2-> A warp's loads and stores are \alert{coalesced} into one
      transaction when its 32 addresses are adjacent. With x on \cd{j}, a
      warp writes 32 adjacent floats of \cd{c}, and reads 32 of \cd{b}.
  \end{itemize}
\end{frame}
\note{No timing, deliberately: Colab GPU timings are noise, and a number that
moves for the wrong reason teaches worse than no number.

The same goes for a: with x on j every lane of a warp reads the same
a[i, k]; with x on i the lanes read a[i, k] 2 KB apart, and every lane reads
the same b[k, j].

Swapping y and x in both mappings is the right-hand picture: each lane is a
row of c apart, one transaction per lane. It runs and matches numpy; the
notebook's check is what rejects it. That is why the exercise put x on the
columns, as Session 1 Part 2's kernel did; in the schedule it is one
attribute per loop.}

% =======================================================================
\begin{frame}{Recap}
  \begin{center}
  \begin{tikzpicture}[
      box/.style={draw=scblue,fill=scwash,rounded corners=2pt,inner sep=5pt,
                  font=\ttfamily\scriptsize,minimum height=2.4em,align=center},
      lbl/.style={font=\scriptsize\itshape,color=scgrey,align=center},
      arr/.style={-{Stealth[length=2mm]},draw=scgrey,thick}]
    \node[box] (op) at (0,0) {linalg.matmul};
    \node[box,draw=scaccent] (forall) at (4.0,0)
      {scf.forall $\times$ 2\\\normalfont\itshape mapped: blocks, threads};
    \node[box] (launch) at (7.9,0)
      {gpu.launch\\\normalfont\itshape 8$\times$4 blocks of 32$\times$32};
    \node[box] (run) at (11.3,0) {JIT + run\\\normalfont\itshape T4 GPU};
    \draw[arr,draw=scaccent] (op) -- node[lbl,above=2pt]{tile} (forall);
    \draw[arr,draw=scaccent] (forall) -- node[lbl,above=2pt]{map} (launch);
    \draw[arr] (launch) -- node[lbl,above=2pt]{GPU\\pipeline} (run);
    \node[lbl] at (4.0,-0.95) {all in the schedule};
    \draw[scgrey,thin] (1.0,-0.7) -- (6.6,-0.7);
  \end{tikzpicture}
  \end{center}
  \vspace{0.3em}
  \begin{itemize}
    \item<2-> The same tiling as Part 1, into \cd{scf.forall} loops, each
      mapped onto the GPU: the outer one to blocks, the nested one to
      threads.
    \item<3-> The schedule states the mapping, so it also decides how memory
      is accessed: x on the columns, and each warp touches adjacent floats.
  \end{itemize}
\end{frame}
\note{End here.}

\end{document}
